\chapter{Motivación y objetivos}

\section{Introducción}

En 2019 OpenAI presentó GPT-2, un modelo de lenguaje basado en la arquitectura \textit{Transformer} \cite{radford2019gpt2}. A diferencia de arquitecturas recurrentes, los \textit{Transformers} permiten procesar secuencias mediante un mecanismo de atención que relaciona directamente los distintos elementos de la entrada. Esta arquitectura fue presentada en \textit{Attention Is All You Need} \cite{vaswani2017attention} y se ha convertido en uno de los componentes fundamentales de los modelos de lenguaje modernos.

Desde entonces, el tamaño de los modelos, la cantidad de datos utilizados durante el entrenamiento y la longitud de las secuencias procesadas han aumentado considerablemente. Esta evolución ha estado acompañada por un incremento igualmente importante de los recursos computacionales necesarios tanto para entrenamiento como para inferencia.

Uno de los principales costes de los modelos basados en \textit{Transformers} se encuentra precisamente en el mecanismo de atención. Para una secuencia de longitud \(N\), la matriz de puntuaciones

$$
QK^T
$$

contiene \(N^2\) elementos. En la implementación convencional de la atención, esta matriz debe materializarse y posteriormente ser procesada por la operación de \textit{softmax}, generando una cantidad considerable de tráfico a memoria.

El coste computacional de la atención es igualmente cuadrático respecto a la longitud de secuencia. Como consecuencia, el incremento de las ventanas de contexto utilizadas por los modelos actuales convierte esta operación en un objetivo especialmente relevante para la optimización.

Las unidades de procesamiento gráfico (GPU) se han convertido en el principal tipo de acelerador utilizado para estas cargas de trabajo. Su elevado paralelismo y la incorporación de unidades especializadas para operaciones matriciales permiten alcanzar un gran rendimiento, pero obtener dicho rendimiento requiere considerar explícitamente aspectos como la jerarquía de memoria, la distribución del trabajo, la sincronización entre hilos y la utilización de los recursos internos de cada multiprocesador.

Paralelamente, han aparecido lenguajes y compiladores especializados para facilitar el desarrollo de kernels de alto rendimiento y automatizar parte de su optimización. Entre ellos se encuentra Triton, utilizado en este trabajo junto con CUDA y con implementaciones de referencia de PyTorch.

\section{Motivación}

FlashAttention constituye un ejemplo especialmente adecuado para estudiar este problema. El algoritmo reorganiza el cálculo de la atención para evitar la materialización completa de las matrices intermedias en memoria global. Para ello, divide las matrices \(Q\), \(K\) y \(V\) en bloques y mantiene los resultados intermedios en niveles más rápidos de la jerarquía de memoria \cite{dao2022flashattention}.

Esta estrategia reduce considerablemente el tráfico hacia memoria global, pero traslada parte de la complejidad al diseño del kernel. Una implementación eficiente debe decidir cómo distribuir los bloques entre los \textit{warps}, qué datos conservar en registros o memoria compartida, cómo realizar las transferencias entre niveles de memoria y cómo alimentar de forma continua las unidades matriciales de la GPU.

La arquitectura Ampere proporciona mecanismos específicamente diseñados para este tipo de cargas de trabajo \cite{nvidiaAmpere}. Entre ellos se encuentran las instrucciones asíncronas \texttt{cp.async}, las cargas colectivas \texttt{ldmatrix} y las instrucciones \texttt{mma.sync} empleadas para ejecutar productos matriciales sobre los Tensor Cores; su semántica se documenta en la ISA PTX de NVIDIA \cite{nvidiaPTX}.

El control directo de estas instrucciones mediante CUDA permite estudiar detalladamente su comportamiento, pero también incrementa considerablemente la complejidad de implementación. Triton, por el contrario, permite expresar gran parte del cálculo utilizando operaciones matriciales de alto nivel y delegar al compilador decisiones como el reparto de los datos, el uso de memoria compartida o la generación de determinadas instrucciones de bajo nivel.

La motivación principal de este trabajo consiste, por tanto, en implementar FlashAttention utilizando ambos enfoques y estudiar las diferencias entre ellos. El interés no reside únicamente en determinar qué versión obtiene un menor tiempo de ejecución, sino en relacionar el rendimiento observado con decisiones concretas de organización del trabajo, movimiento de datos y utilización de los recursos de la GPU.

La comparación se completa utilizando como referencia \path{scaled_dot_product_attention} (SDPA) de PyTorch \cite{pytorchSDPA} y su backend basado en cuDNN \cite{cudnnAttention}. Estas rutas permiten situar el rendimiento obtenido por los kernels desarrollados dentro de un contexto realista.

\section{Objetivos}

El objetivo general del trabajo consiste en estudiar la implementación y optimización de FlashAttention sobre GPUs NVIDIA, comparando una implementación propia en CUDA, una implementación desarrollada con Triton y las rutas de referencia disponibles mediante PyTorch SDPA y cuDNN.

A partir de este objetivo general se establecen los siguientes objetivos específicos:

\begin{itemize}

\item Estudiar el funcionamiento del mecanismo de atención utilizado por los modelos \textit{Transformer} y analizar los principales costes computacionales y de memoria asociados a su implementación convencional.

\item Estudiar el algoritmo FlashAttention y, en particular, la utilización del \textit{softmax online} para evitar la materialización completa de la matriz de atención.

\item Desarrollar una implementación de FlashAttention directamente en CUDA para GPUs NVIDIA de arquitectura Ampere y dimensión de cabeza \(d=128\).

\item Utilizar explícitamente los Tensor Cores mediante instrucciones MMA con operandos BF16 y acumulación FP32.

\item Implementar un \textit{pipeline} de transferencia entre memoria global y memoria compartida mediante instrucciones \texttt{cp.async} y múltiples buffers.

\item Utilizar instrucciones \texttt{ldmatrix} para transferir los fragmentos necesarios desde memoria compartida al fichero de registros.

\item Diseñar un patrón de \textit{swizzling} para la memoria compartida que permita evitar conflictos de banco durante las cargas matriciales.

\item Optimizar los accesos a memoria global de \(Q\), \(K\), \(V\) y \(O\), buscando que las transferencias realizadas por cada \textit{warp} sean completamente coalescentes.

\item Desarrollar una segunda implementación equivalente utilizando Triton y explorar sus parámetros de ejecución.

\item Utilizar el mecanismo de \textit{autotuning} de Triton para estudiar diferentes configuraciones de tamaño de bloque, número de \textit{warps} y etapas del \textit{pipeline}.

\item Implementar un test funcional de referencia que compruebe la correctitud numérica de los kernels mediante \texttt{torch.allclose} frente a SDPA.

\item Desarrollar tests no funcionales destinados a medir el tiempo de ejecución y estudiar la escalabilidad de las distintas implementaciones al incrementar la longitud de secuencia.

\item Analizar los kernels utilizando NVIDIA Nsight Compute, prestando especial atención a métricas relacionadas con ocupación, utilización de registros y memoria compartida, accesos no coalescentes, conflictos de banco, utilización de los Tensor Cores y causas de espera de los \textit{warps}.

\item Comparar el rendimiento de las implementaciones CUDA y Triton con SDPA y con el backend basado en cuDNN.

\item Identificar las decisiones de diseño que limitan el rendimiento de la implementación CUDA y proponer posibles líneas de mejora.

\end{itemize}

\section{Originalidad del TFG}

La principal aportación del trabajo reside en el desarrollo desde cero de una implementación de FlashAttention directamente sobre CUDA utilizando primitivas específicas de la arquitectura Ampere.

En lugar de limitar la comparación a tiempos de ejecución de implementaciones existentes, se estudia el algoritmo desde el nivel de las instrucciones ejecutadas por la GPU. Esto incluye el reparto explícito del trabajo entre \textit{warps}, la utilización del fichero de registros y de la memoria compartida, la construcción de un \textit{pipeline} mediante transferencias asíncronas, el acceso a Tensor Cores y la reorganización física de los datos mediante \textit{swizzling}.

La implementación CUDA se compara con una versión funcionalmente equivalente desarrollada mediante Triton y con las implementaciones de referencia de PyTorch. La aportación no se limita a ordenar estas rutas por tiempo de ejecución, sino que busca relacionar sus diferencias con métricas hardware obtenidas mediante perfilado.

Finalmente, el análisis mediante Nsight Compute permite relacionar las diferencias de rendimiento con características concretas de los kernels, en lugar de limitar la evaluación a una comparación puramente temporal.

\section{Planificación}

El desarrollo del proyecto se ha realizado de forma incremental. En una primera fase se estudió el algoritmo de atención y el funcionamiento de FlashAttention, prestando especial atención al \textit{softmax online} y a la descomposición por bloques necesaria para evitar la materialización de la matriz completa de atención.

Posteriormente se desarrolló una primera implementación CUDA funcional. Sobre esta versión se fueron incorporando progresivamente las distintas optimizaciones de bajo nivel.

Entre las principales etapas del desarrollo se encuentran:

\begin{enumerate}

\item implementación del producto \(QK^T\) utilizando Tensor Cores;

\item implementación del \textit{softmax online};

\item acumulación del producto \(PV\) manteniendo la salida en FP32;

\item introducción de transferencias asíncronas mediante \texttt{cp.async};

\item utilización de múltiples buffers en memoria compartida;

\item incorporación de las cargas mediante \texttt{ldmatrix};

\item diseño del patrón de \textit{swizzling} de memoria compartida;

\item reorganización de las cargas globales de \(Q\), \(K\) y \(V\) para eliminar accesos no coalescentes;

\item reorganización de la escritura de \(O\) para conseguir stores coalescentes;

\item análisis del consumo de registros, memoria compartida y ocupación;

\item desarrollo de la implementación Triton y exploración de sus parámetros de ejecución;

\item desarrollo del test funcional y de los tests de rendimiento;

\item evaluación de las implementaciones mediante NVIDIA Nsight Compute;

\item comparación final con SDPA y cuDNN.

\end{enumerate}

El proceso de optimización ha sido iterativo. Después de cada modificación se verificaba en primer lugar la correctitud del kernel mediante el test funcional y, posteriormente, se analizaba su efecto sobre el rendimiento y las métricas obtenidas mediante el \textit{profiler}. Esta metodología permitió evitar que una optimización de rendimiento introdujera errores silenciosos en el resultado.

\section{Organización del documento}

El resto de la memoria se organiza de la siguiente forma.

En primer lugar, se presentan los conceptos necesarios para comprender el mecanismo de atención, la arquitectura \textit{Transformer}, FlashAttention y las características de las GPUs NVIDIA relevantes para su implementación.

Posteriormente se describe la metodología de desarrollo utilizada. Se presenta el algoritmo FlashAttention empleado como referencia y se detallan las implementaciones CUDA y Triton, incluyendo la distribución del trabajo, la utilización de la jerarquía de memoria, los Tensor Cores y las principales optimizaciones aplicadas.

A continuación se describe el entorno de prueba utilizado para realizar las mediciones, incluyendo las características de la GPU NVIDIA A100, el entorno software y las versiones de las herramientas empleadas.

El siguiente capítulo presenta el test funcional utilizado para verificar la correctitud y los tests no funcionales empleados en las medidas temporales y en el \textit{profiling}.

Posteriormente se realiza el análisis de los resultados obtenidos. Se comparan las implementaciones CUDA, Triton, SDPA y cuDNN mediante sus tiempos de ejecución y mediante las métricas obtenidas con NVIDIA Nsight Compute. Se estudian aspectos como la ocupación, el número de registros, el consumo de memoria compartida, los conflictos de banco, los accesos coalescentes y el número total de sectores de memoria solicitados.

Finalmente, se presentan las conclusiones del trabajo, las principales lecciones obtenidas durante el desarrollo y las posibles líneas de investigación y optimización futuras.
